\documentclass[a4paper,12pt]{article}

\usepackage[T2A]{fontenc}
\usepackage[utf8]{inputenc}
\usepackage[russian]{babel}
\usepackage{amsmath,amssymb}
\usepackage{paratype}
\renewcommand{\familydefault}{\sfdefault}
\usepackage{icomma}
\usepackage[a4paper,top=19mm,bottom=19mm,left=20mm,right=20mm,headheight=15pt]{geometry}
\usepackage{microtype}
\usepackage{tikz}
\usetikzlibrary{calc,arrows.meta,positioning,shapes.geometric}
\usepackage{pgfplots}
\pgfplotsset{compat=1.18}
\usepackage{tabularx,booktabs}
\usepackage{enumitem}
\usepackage{hyperref}
\usepackage{parskip}
\usepackage{fancyhdr}
\usepackage{setspace}
\usepackage{xcolor}

\definecolor{accent}{HTML}{008080}
\definecolor{darkaccent}{HTML}{004D4D}
\definecolor{textgray}{HTML}{333333}
\definecolor{softgray}{HTML}{F2F4F4}
\definecolor{lightaccent}{HTML}{E4F1F1}

\hypersetup{
  colorlinks=true,
  linkcolor=darkaccent,
  urlcolor=darkaccent,
  pdftitle={Чек-лист: парабола, гипербола и модуль},
  pdfauthor={Артём Александрович Лёвин}
}

\onehalfspacing
\setlength{\parindent}{0pt}
\setlength{\parskip}{0.55em}
\renewcommand{\arraystretch}{1.25}

\setlist[itemize]{
  leftmargin=1.55em,
  itemsep=0.22em,
  topsep=0.25em
}
\setlist[enumerate]{
  leftmargin=1.85em,
  itemsep=0.35em,
  topsep=0.3em
}

\pagestyle{fancy}
\fancyhf{}
\fancyhead[L]{\small\color{textgray}Ярослав Верещагин \textbullet{} ОГЭ}
\fancyhead[R]{\small\color{textgray}Парабола, гипербола, модуль}
\fancyfoot[C]{\small\color{textgray}\thepage}
\renewcommand{\headrulewidth}{0.3pt}
\renewcommand{\footrulewidth}{0pt}

\newcommand{\topic}[1]{%
  \vspace{0.25em}
  {\Large\bfseries\color{darkaccent}#1}\par
  \vspace{0.1em}
  {\color{accent}\rule{\linewidth}{0.7pt}}\par
  \vspace{0.35em}
}

\newcommand{\subtopic}[1]{%
  \vspace{0.35em}
  {\large\bfseries\color{darkaccent}#1}\par
  \vspace{0.1em}
}

\newcommand{\key}[1]{\textbf{\color{darkaccent}#1}}

\tikzset{
  flowstart/.style={
    ellipse,
    draw=accent,
    very thick,
    minimum width=40mm,
    minimum height=10mm,
    align=center,
    fill=lightaccent,
    font=\small
  },
  flowaction/.style={
    rounded corners=3pt,
    draw=accent,
    thick,
    minimum width=52mm,
    minimum height=12mm,
    text width=48mm,
    align=center,
    fill=white,
    font=\small
  },
  flowdecision/.style={
    diamond,
    aspect=2.2,
    draw=darkaccent,
    thick,
    text width=34mm,
    align=center,
    inner sep=1.5pt,
    fill=softgray,
    font=\small
  },
  flowend/.style={
    ellipse,
    draw=darkaccent,
    very thick,
    minimum width=42mm,
    minimum height=10mm,
    align=center,
    fill=lightaccent,
    font=\small
  },
  flowarrow/.style={
    -{Stealth[length=3mm,width=2mm]},
    thick,
    draw=darkaccent
  }
}

\begin{document}

% =================================================
% Титульный лист
% =================================================

\begin{titlepage}
\thispagestyle{empty}
\centering

\vspace*{22mm}

{\color{accent}\rule{0.72\linewidth}{1.2pt}}\\[13mm]

{\fontsize{28}{33}\selectfont
 \bfseries\color{darkaccent}
 Чек-лист}\\[5mm]

{\fontsize{21}{26}\selectfont
 \bfseries
 Парабола, гипербола и модуль}\\[8mm]

{\large
Графики функций: чтение, построение и преобразования}\\[17mm]

{\large\bfseries
Ярослав Верещагин}\\[2mm]

{\normalsize
Подготовка к ОГЭ по математике}\\[19mm]

{\normalsize
Артём Александрович Лёвин}\\[1mm]

{\normalsize
эксклюзивно для Ярослава Верещагина}\\[12mm]

{\large 24 сентября 2026 г.}

\vfill

\begin{tikzpicture}[remember picture,overlay]
  \draw[
    accent,
    line width=1.15pt
  ]
    ($(current page.south west)+(18mm,20mm)$)
    .. controls
    ($(current page.south west)+(64mm,33mm)$)
    and
    ($(current page.south east)+(-72mm,8mm)$)
    ..
    ($(current page.south east)+(-18mm,22mm)$);

  \draw[
    accent!45,
    line width=0.55pt
  ]
    ($(current page.south west)+(20mm,15mm)$)
    .. controls
    ($(current page.south west)+(84mm,28mm)$)
    and
    ($(current page.south east)+(-91mm,3mm)$)
    ..
    ($(current page.south east)+(-20mm,17mm)$);
\end{tikzpicture}

\end{titlepage}

% =================================================
% Парабола
% =================================================

\topic{1. Парабола: читаем коэффициенты по графику}

Рассматриваем квадратичную функцию
\[
  y=ax^2+bx+c,
  \qquad a\ne0.
\]

\subtopic{Что означает каждый коэффициент}

\begin{itemize}
  \item
  \key{Коэффициент $a$} задаёт направление ветвей:
  \[
    a>0 \Rightarrow \text{ветви направлены вверх},
    \qquad
    a<0 \Rightarrow \text{ветви направлены вниз}.
  \]

  \item
  \key{Коэффициент $c$} равен значению функции при $x=0$:
  \[
    c=f(0).
  \]
  Поэтому парабола пересекает ось $Oy$ в точке
  \[
    (0;c).
  \]
  Если эта точка выше оси $Ox$, то $c>0$; ниже --- $c<0$;
  если парабола проходит через начало координат, то $c=0$.

  \item
  \key{Абсцисса вершины} вычисляется по формуле
  \[
    x_{\text{в}}=-\frac{b}{2a}.
  \]
  Отсюда
  \[
    b=-2ax_{\text{в}}.
  \]
  Эта запись особенно удобна, когда по рисунку нужно определить только знак $b$.
\end{itemize}

\begin{center}
\begin{tabular}{ccc}
\toprule
Знак $a$ & Знак $x_{\text{в}}$ & Знак $b=-2ax_{\text{в}}$ \\
\midrule
$+$ & $+$ & $-$ \\
$+$ & $-$ & $+$ \\
$-$ & $+$ & $+$ \\
$-$ & $-$ & $-$ \\
\bottomrule
\end{tabular}
\end{center}

Если вершина лежит на оси $Oy$, то
\[
  x_{\text{в}}=0
  \quad\Rightarrow\quad
  b=0.
\]

\subtopic{Пример}

Пусть ветви параболы направлены вверх, точка пересечения с осью $Oy$
лежит ниже оси $Ox$, а вершина находится слева от оси $Oy$.
Тогда
\[
  a>0,
  \qquad
  c<0,
  \qquad
  x_{\text{в}}<0.
\]
Следовательно,
\[
  b=-2ax_{\text{в}}>0.
\]

\begin{figure}[ht]
\centering
\begin{tikzpicture}
\begin{axis}[
  width=0.76\linewidth,
  height=7cm,
  axis lines=middle,
  xmin=-5,
  xmax=4,
  ymin=-5,
  ymax=6,
  xtick={-4,-3,-2,-1,1,2,3},
  ytick={-4,-2,2,4},
  xlabel={$x$},
  ylabel={$y$},
  axis equal image,
  unit vector ratio=1 1,
  samples=240,
  clip=false,
  grid=both,
  minor tick num=0,
  grid style={line width=.15pt,draw=gray!25},
  every axis plot/.append style={very thick}
]
\addplot[accent,domain=-4.2:2.6] {0.7*(x+1.5)^2-3};
\addplot[only marks,mark=*,mark size=2.1pt,darkaccent]
coordinates {(-1.5,-3) (0,-1.425)};
\node[anchor=north east,text=darkaccent] at (axis cs:-1.5,-3) {$V$};
\node[anchor=west,text=darkaccent] at (axis cs:0,-1.425) {$(0;c)$};
\end{axis}
\end{tikzpicture}
\end{figure}

\key{Самопроверка.}
В формуле $b=-2ax_{\text{в}}$ не теряйте внешний минус.
Сначала определите знаки $a$ и $x_{\text{в}}$, затем отдельно выполните анализ знаков.

% =================================================
% Блок-схема параболы
% =================================================

\clearpage
\thispagestyle{fancy}

\begin{center}
{\Large\bfseries\color{darkaccent}
Алгоритм: определить знаки $a$, $b$, $c$ по графику параболы}
\end{center}

\vspace{7mm}

\begin{center}
\begin{tikzpicture}[node distance=11mm and 10mm]

\node[flowstart] (start) {Дан график\\$y=ax^2+bx+c$};

\node[flowdecision,below=of start] (a) {Куда направлены ветви?};

\node[flowaction,below left=of a] (apos) {$a>0$\\ветви вверх};
\node[flowaction,below right=of a] (aneg) {$a<0$\\ветви вниз};

\node[flowaction,below=27mm of a] (c)
{Найдите пересечение с $Oy$:\\его ордината равна $c$};

\node[flowaction,below=of c] (xv)
{Определите знак $x_{\text{в}}$\\по положению вершины};

\node[flowaction,below=of xv] (b)
{Определите знак $b$ по формуле\\$b=-2ax_{\text{в}}$};

\node[flowend,below=of b] (end)
{Запишите знаки $a$, $b$, $c$};

\draw[flowarrow] (start) -- (a);
\draw[flowarrow] (a) -| node[pos=0.22,above left]{\small вверх} (apos);
\draw[flowarrow] (a) -| node[pos=0.22,above right]{\small вниз} (aneg);
\draw[flowarrow] (apos) |- (c);
\draw[flowarrow] (aneg) |- (c);
\draw[flowarrow] (c) -- (xv);
\draw[flowarrow] (xv) -- (b);
\draw[flowarrow] (b) -- (end);

\end{tikzpicture}
\end{center}

\clearpage

% =================================================
% Гипербола
% =================================================

\topic{2. Гипербола $y=\dfrac{k}{x}$}

Рассматриваем функцию
\[
  y=\frac{k}{x},
  \qquad k\ne0.
\]

\subtopic{ОДЗ и основные свойства}

\begin{itemize}
  \item
  Делить на ноль нельзя, поэтому
  \[
    x\ne0.
  \]
  Область определения:
  \[
    (-\infty;0)\cup(0;+\infty).
  \]

  \item
  Значение $y=0$ для такой функции невозможно, поэтому
  \[
    y\ne0.
  \]

  \item
  Прямые
  \[
    x=0
    \qquad\text{и}\qquad
    y=0
  \]
  являются асимптотами графика.

  \item
  Если $k>0$, ветви расположены в I и III четвертях.

  \item
  Если $k<0$, ветви расположены во II и IV четвертях.
\end{itemize}

\subtopic{Как влияет величина $|k|$}

При одном и том же ненулевом $x$
\[
  |y|=\frac{|k|}{|x|}.
\]
Поэтому при увеличении $|k|$ соответствующая точка графика имеет большее $|y|$;
при уменьшении $|k|$ --- меньшее $|y|$.

Например, при $x=1$
\[
  \frac{2}{1}=2,
  \qquad
  \frac{1}{1}=1,
  \qquad
  \frac{0{,}5}{1}=0{,}5.
\]

\subtopic{Коэффициент в знаменателе}

Формулу полезно привести к стандартному виду $k/x$:
\[
  \frac{1}{2x}
  =
  \frac{1/2}{x}
  =
  \frac{0{,}5}{x}.
\]
Следовательно, здесь
\[
  k=0{,}5.
\]

\subtopic{Построение по точкам}

Ноль брать нельзя. Удобно выбирать значения $x$, симметричные относительно нуля,
например
\[
  -4,-2,-1,1,2,4.
\]
Для каждого значения вычисляем
\[
  y=\frac{k}{x}.
\]

\begin{figure}[ht]
\centering
\begin{tikzpicture}
\begin{axis}[
  width=0.80\linewidth,
  height=7cm,
  axis lines=middle,
  xmin=-5,
  xmax=5,
  ymin=-5,
  ymax=5,
  xlabel={$x$},
  ylabel={$y$},
  axis equal image,
  unit vector ratio=1 1,
  samples=260,
  clip=true,
  grid=both,
  grid style={line width=.15pt,draw=gray!25},
  legend style={
    at={(0.5,-0.16)},
    anchor=north,
    legend columns=2,
    draw=none,
    font=\small
  }
]
\addplot[accent,very thick,domain=-5:-0.18,forget plot] {1/x};
\addplot[accent,very thick,domain=0.18:5] {1/x};
\addlegendentry{$y=1/x$}

\addplot[darkaccent,very thick,dashed,domain=-5:-0.18,forget plot] {2/x};
\addplot[darkaccent,very thick,dashed,domain=0.18:5] {2/x};
\addlegendentry{$y=2/x$}

\addplot[gray!70!black,very thick,dotted,domain=-5:-0.18,forget plot] {0.5/x};
\addplot[gray!70!black,very thick,dotted,domain=0.18:5] {0.5/x};
\addlegendentry{$y=0{,}5/x$}

\addplot[black,thick,dashdotted,domain=-5:-0.18,forget plot] {-1/x};
\addplot[black,thick,dashdotted,domain=0.18:5] {-1/x};
\addlegendentry{$y=-1/x$}
\end{axis}
\end{tikzpicture}
\end{figure}

\key{Типичные ошибки:}
\begin{itemize}[itemsep=0.08em,topsep=0.12em]
  \item подставлять $x=0$;
  \item соединять две ветви через ось $Oy$;
  \item путать четверти при $k>0$ и $k<0$;
  \item считать, что в $y=1/(2x)$ коэффициент $k$ равен $2$.
\end{itemize}

% =================================================
% Модуль
% =================================================

\clearpage
\topic{3. Модуль, наложенный на всю правую часть}

Рассматриваем функцию
\[
  y=|f(x)|.
\]

Здесь модуль охватывает \emph{всю} правую часть. Именно для такого случая действует правило отражения относительно оси $Ox$.

\[
  |f(x)|=
  \begin{cases}
    f(x), & f(x)\ge0,\\
    -f(x), & f(x)<0.
  \end{cases}
\]

Отсюда следует геометрическое правило:

\begin{itemize}
  \item часть графика $y=f(x)$, лежащая на оси $Ox$ или выше неё, остаётся на месте;
  \item часть графика, лежащая ниже оси $Ox$, отражается симметрично относительно оси $Ox$ вверх;
  \item у графика $y=|f(x)|$ всегда выполняется $y\ge0$.
\end{itemize}

\subtopic{Где искать точки смены правила}

Нужно решить уравнение
\[
  f(x)=0.
\]
Именно в этих точках исходный график пересекает ось $Ox$,
а знак $f(x)$ может измениться.

Для линейного выражения под модулем корень даёт вершину «галочки».

Например,
\[
  y=|2x-1|.
\]
Точка излома находится из условия
\[
  2x-1=0,
\]
поэтому
\[
  x=\frac12.
\]
Вершина графика имеет координаты
\[
  \left(\frac12;0\right).
\]

\begin{figure}[ht]
\centering
\begin{tikzpicture}
\begin{axis}[
  width=0.82\linewidth,
  height=7.2cm,
  axis lines=middle,
  xmin=-3,
  xmax=4,
  ymin=-4,
  ymax=6,
  xlabel={$x$},
  ylabel={$y$},
  axis equal image,
  unit vector ratio=1 1,
  samples=220,
  clip=false,
  grid=both,
  grid style={line width=.15pt,draw=gray!25},
  legend style={
    at={(0.5,-0.16)},
    anchor=north,
    legend columns=2,
    draw=none,
    font=\small
  }
]
\addplot[gray!65,thick,dashed,domain=-3:4] {2*x-1};
\addlegendentry{$y=2x-1$}
\addplot[accent,very thick,domain=-3:4] {abs(2*x-1)};
\addlegendentry{$y=|2x-1|$}
\addplot[only marks,mark=*,mark size=2.1pt,darkaccent]
coordinates {(0.5,0)};
\node[anchor=north east,text=darkaccent] at (axis cs:0.5,0) {$\left(\frac12;0\right)$};
\end{axis}
\end{tikzpicture}
\end{figure}

\subtopic{То же правило для параболы}

Для
\[
  y=|x^2-x-6|
\]
сначала рассматриваем
\[
  y=x^2-x-6=(x-3)(x+2).
\]
Корни исходной параболы:
\[
  x=-2,
  \qquad
  x=3.
\]
Поскольку ветви исходной параболы направлены вверх, она лежит ниже оси $Ox$
на промежутке
\[
  (-2;3).
\]
Именно этот участок при переходе к модулю отражается вверх.

\subtopic{То же правило для гиперболы}

Для
\[
  y=\left|\frac1x\right|
\]
сохраняется ограничение
\[
  x\ne0.
\]
Ветвь исходной гиперболы $y=1/x$, лежащая в III четверти,
отражается во II четверть. Итоговый график целиком расположен выше оси $Ox$.

\key{Самопроверка.}
После построения $y=|f(x)|$ на рисунке не должно остаться ни одной точки с $y<0$.

% =================================================
% Блок-схема модуля
% =================================================

\clearpage
\thispagestyle{fancy}

\begin{center}
{\Large\bfseries\color{darkaccent}
Алгоритм: построить график $y=|f(x)|$}
\end{center}

\vspace{8mm}

\begin{center}
\begin{tikzpicture}[node distance=12mm and 10mm]

\node[flowstart] (start) {Дана функция\\$y=|f(x)|$};

\node[flowaction,below=of start] (base)
{Постройте исходный график\\$y=f(x)$};

\node[flowaction,below=of base] (roots)
{Найдите точки, где\\$f(x)=0$};

\node[flowdecision,below=of roots] (sign)
{Есть участки графика\\ниже оси $Ox$?};

\node[flowaction,below left=of sign] (reflect)
{Отразите все такие участки\\симметрично относительно $Ox$};

\node[flowaction,below right=of sign] (keep)
{Оставьте исходный график\\без изменений};

\node[flowaction,below=28mm of sign] (check)
{Проверьте итог:\\все точки должны удовлетворять $y\ge0$};

\node[flowend,below=of check] (end)
{График $y=|f(x)|$ построен};

\draw[flowarrow] (start) -- (base);
\draw[flowarrow] (base) -- (roots);
\draw[flowarrow] (roots) -- (sign);
\draw[flowarrow] (sign) -| node[pos=0.22,above left]{\small да} (reflect);
\draw[flowarrow] (sign) -| node[pos=0.22,above right]{\small нет} (keep);
\draw[flowarrow] (reflect) |- (check);
\draw[flowarrow] (keep) |- (check);
\draw[flowarrow] (check) -- (end);

\end{tikzpicture}
\end{center}

\clearpage

% =================================================
% Итоги и домашняя подготовка
% =================================================

\topic{4. Что нужно уметь после занятия}

\begin{enumerate}
  \item
  По графику параболы $y=ax^2+bx+c$ определять знак $a$,
  знак $c$, положение вершины и затем знак $b$ по формуле
  \[
    b=-2ax_{\text{в}}.
  \]

  \item
  Для гиперболы $y=k/x$ помнить ограничение $x\ne0$,
  распознавать четверти по знаку $k$ и учитывать влияние величины $|k|$.

  \item
  Строить гиперболу по таблице значений, не используя $x=0$.

  \item
  Для $y=|f(x)|$ оставлять верхнюю часть исходного графика на месте,
  а нижнюю отражать относительно оси $Ox$.

  \item
  Для линейного выражения под модулем находить точку излома из уравнения
  \[
    f(x)=0.
  \]
\end{enumerate}

\subtopic{Домашняя подготовка к текстовым задачам ОГЭ}

Для каждого из заданий 1--5:
\begin{enumerate}
  \item внимательно прочитайте весь текст задачи;
  \item выпишите все введённые величины и обозначения;
  \item оформите краткую запись «Дано»;
  \item отдельно запишите, что требуется найти;
  \item решение пока не выполняйте.
\end{enumerate}

Цель такой подготовки --- заранее разобраться в условии и обозначениях,
чтобы на занятии сосредоточиться на построении решения.

% =================================================
% Тренировочный блок
% =================================================

\clearpage

\topic{Тренировочный блок --- Путь А}

\begingroup
\small
\setstretch{1.18}
\setlist[enumerate]{leftmargin=1.8em,itemsep=0.2em,topsep=0.18em}

\textbf{5 заданий.}
Задания расположены от базового чтения графика к комбинированному применению правил.

\begin{enumerate}
  \item
  Парабола
  \[
    y=ax^2+bx+c
  \]
  имеет ветви вверх, пересекает ось $Oy$ ниже оси $Ox$,
  а её вершина лежит слева от оси $Oy$.

  Определите знаки
  \[
    a,\qquad b,\qquad c.
  \]

  \item
  Установите соответствие между функциями и описаниями их графиков:
  \[
    \text{A) } y=\frac{3}{x},
    \qquad
    \text{Б) } y=-\frac{2}{x},
    \qquad
    \text{В) } y=\frac{1}{4x}.
  \]

  \begin{enumerate}[label=\arabic*),itemsep=0.1em]
    \item ветви расположены во II и IV четвертях;
    \item ветви расположены в I и III четвертях, причём при $|x|=1$ имеем $|y|>1$;
    \item ветви расположены в I и III четвертях, причём при $|x|=1$ имеем $|y|<1$.
  \end{enumerate}

  \item
  Для функции
  \[
    y=\frac{2}{x}
  \]
  составьте таблицу значений при
  \[
    x\in\{-4,-2,-1,1,2,4\},
  \]
  постройте график и укажите его асимптоты.

  \item
  Постройте график
  \[
    y=|x-2|.
  \]
  Укажите координаты точки излома и определите,
  какая часть прямой $y=x-2$ отражается относительно оси $Ox$.

  \item
  Для функции
  \[
    y=|x^2-4x+3|
  \]
  определите без подробного построения:
  \begin{enumerate}[label=\alph*),itemsep=0.1em]
    \item корни подмодульного выражения;
    \item абсциссу вершины исходной параболы;
    \item промежуток, на котором исходная парабола расположена ниже оси $Ox$;
    \item участок, который при переходе к модулю отражается вверх.
  \end{enumerate}
\end{enumerate}
\endgroup

% =================================================
% Ответы
% =================================================

\clearpage

\topic{Ответы}

\begin{table}[ht]
\centering
\begin{tabularx}{\linewidth}{@{}>{\bfseries}c X@{}}
\toprule
№ & Ответ \\
\midrule
1 &
$\displaystyle a>0,\quad b>0,\quad c<0.$
\\[2mm]

2 &
A--2, Б--1, В--3.
\\[2mm]

3 &
При $x=-4,-2,-1,1,2,4$ соответственно
\[
  y=-\frac12,-1,-2,2,1,\frac12.
\]
Асимптоты:
\[
  x=0,\qquad y=0.
\]
\\[2mm]

4 &
Точка излома:
\[
  (2;0).
\]
Участок прямой $y=x-2$, соответствующий $x<2$, лежит ниже оси $Ox$
и отражается вверх.
\\[2mm]

5 &
\[
  x^2-4x+3=(x-1)(x-3),
\]
поэтому корни: $1$ и $3$;
\[
  x_{\text{в}}=2.
\]
Исходная парабола ниже оси $Ox$ на промежутке $(1;3)$,
и именно этот участок отражается вверх.
\\
\bottomrule
\end{tabularx}
\end{table}

\vfill

{\small\color{textgray}
Короткая самопроверка перед завершением работы:
для параболы контролируйте минус в $b=-2ax_{\text{в}}$;
для гиперболы --- запрет $x=0$;
для $y=|f(x)|$ --- отсутствие точек с $y<0$.
}

\end{document}